% ═══════════════════════════════════════════════════════════════════════════
% QR-CODE: Lernpfad Ohmsches Gesetz (Kl. 8, DS 02)
% Großer QR-Code auf einer A4-Seite (Beamer-Projektion)
% ═══════════════════════════════════════════════════════════════════════════

\documentclass[a4paper,11pt]{article}

\usepackage[utf8]{inputenc}
\usepackage[T1]{fontenc}
\usepackage{helvet}
\renewcommand{\familydefault}{\sfdefault}

\usepackage[margin=20mm]{geometry}
\usepackage{tikz}
\usepackage{xcolor}
\usepackage{qrcode}
\usepackage{hyperref}

% Hyperref-Einstellungen
\hypersetup{
    colorlinks=false,
    pdfborder={0 0 0}
}

% ═══════════════════════════════════════════════════════════════════════════
% KONFIGURATION
% ═══════════════════════════════════════════════════════════════════════════

\definecolor{physik}{HTML}{2c5aa0}
\colorlet{fachfarbe}{physik}

\newcommand{\lptitel}{Ohmsches Gesetz}
\newcommand{\lpurl}{https://devbox42.github.io/sciencesim/physik/kl08/02-ohm/LP-02a-ohmsches-gesetz.html}
\newcommand{\lpurlkurz}{devbox42.github.io/.../LP-02a-ohmsches-gesetz.html}

% ═══════════════════════════════════════════════════════════════════════════
% DOKUMENT
% ═══════════════════════════════════════════════════════════════════════════

\pagestyle{empty}

\begin{document}

\begin{center}

% Titel
\vspace*{10mm}
{\fontsize{28}{34}\selectfont\bfseries\color{fachfarbe}Lernpfad: \lptitel}

\vspace{15mm}

% QR-Code mit Rahmen (KLICKBAR)
\href{\lpurl}{\fcolorbox{fachfarbe}{white}{\qrcode[height=100mm]{\lpurl}}}

\vspace{12mm}

% URL (klickbar)
{\large\color{gray}\href{\lpurl}{\lpurlkurz}}

\vspace{8mm}

% Hinweis
{\Large\color{gray}Scanne den QR-Code mit deinem Gerät}

\vspace{10mm}

% Info-Box
\fcolorbox{fachfarbe}{fachfarbe!10}{%
    \parbox{0.7\textwidth}{%
        \centering
        \vspace{3mm}
        {\large Bearbeite den Lernpfad vollständig.}\\[2mm]
        Nutze das Arbeitsblatt parallel zum Lernpfad.
        \vspace{3mm}
    }%
}

\end{center}

\end{document}
